\section{Realizaciones sectoriales, composición material y espectros de colecciones}
\label{vi:sec:familias-compuestos}

El atlas de la sección~\ref{vi:sec:atlas} y el lector de la
sección~\ref{vi:sec:operador-masa} tienen funciones complementarias. El
primero clasifica soportes, profundidades y órbitas; el segundo conserva la
cronología de cada trayectoria y determina su operador de masa. Su aplicación
a una colección requiere mantener ambas estructuras. Una igualdad entre
valores de masa no identifica rutas; una igualdad entre celdas tampoco
identifica historias, representaciones o estados compuestos.

En esta sección se desarrollan las realizaciones sectoriales de esa
construcción. Las constantes internas, el par angular, la escala absoluta
y los lectores bidireccionales son los ya construidos a partir de
APP--TRIT--TPK. La representación física precisa qué información del
estado se observa; la comparación metrológica se efectúa después de fijar
esa representación. La exposición conserva, en particular, la diferencia
entre un operador de fibra, una evaluación exacta para una firma declarada
y una línea espectral individualizada por un acoplamiento.

\subsection{Descriptores de colección y criterio espectral de individualización}
\label{vi:subsec:descriptor-familia}

Para un estado elemental escribimos
\[
 X=(\mathfrak c,\rho,g,f,q,\epsilon,\varsigma,\mathfrak a).
\]
Aquí \(\mathfrak c\) es la clase sectorial, \(\rho\) la representación,
\(g\) la generación física, \(f\) el sabor, \(q\) la carga,
\(\epsilon\) la orientación de hoja, \(\varsigma\) la paridad estadística
y \(\mathfrak a\) los datos del acoplamiento. La profundidad combinatoria
del atlas y la generación física siguen siendo coordenadas de tipos
distintos. El sabor distingue componentes de una representación de
colección; el color distingue componentes de la representación cromática;
la hoja y su orientación conservan la procedencia operatoria de la
trayectoria. Ninguna de esas etiquetas sustituye al registro de eventos.

Sea \(\mathcal P_{\rm adm}\) el conjunto de descriptores físicos admisibles,
sea \(\mathcal R_{\rm adm}\) el registro de rutas enriquecidas producido por
el atlas y sea \(\mathcal D\) el espacio de descriptores de incidencia que
conservan clase sectorial, representación, profundidad, carga, color,
orientación y datos de acoplamiento. Definimos
\[
 d_{\mathcal P}:\mathcal P_{\rm adm}\longrightarrow\mathcal D,
 \qquad
 d_{\mathcal R}:\mathcal R_{\rm adm}\longrightarrow\mathcal D,
\]
donde \(d_{\mathcal P}(X)\) expresa los requisitos estructurales del estado
\(X\) y \(d_{\mathcal R}(\gamma)\) expresa las mismas coordenadas conservadas
por la ruta \(\gamma\). La fibra, órbita o multisección producida por el
descriptor es el producto fibrado
\begin{equation}
 F_X=d_{\mathcal R}^{-1}\!\bigl(d_{\mathcal P}(X)\bigr)
 =\mathcal R_{\rm adm}\times_{\mathcal D}\{d_{\mathcal P}(X)\}.
 \label{vi:eq:fibra-descriptor}
\end{equation}
Para los descriptores realizables del artículo, \(F_X\) es finita y no vacía.
En
\(\mathcal H_X=\ell^2(F_X)\otimes V_{\rho_X}\) actúan los dos operadores
\(M_X^{\beta,D}\) y \(M_X^{\beta,\Omega}\) de la
sección~\ref{vi:sec:operador-masa}. Como son diagonales en la base de
rutas, su medida espectral conjunta es
\begin{equation}
 E_X(B)=\sum_{\gamma\in F_X}
 \mathbf1_B\bigl(m_D(\gamma),m_\Omega(\gamma)\bigr)
 |\gamma\rangle\langle\gamma|\otimes I_{V_{\rho_X}},
 \qquad B\subset\mathbb R_{>0}^2.
 \label{vi:eq:familias-pvm}
\end{equation}
Conservar esta medida permite distinguir dos fibras aunque sus medias
coincidan. Un funcional escalar de los dos lectores se aplica únicamente
cuando forma parte de la representación declarada. Si el acoplamiento
conserva varios valores, la salida es un espectro con multiplicidades.

\begin{proposition}[Criterio de línea espectral]
\label{vi:prop:familias-linea}
Sea \(M\) un operador autoadjunto sobre una fibra finita y \(P\) un
proyector de preparación. El canal expresa exclusivamente el valor \(m\)
si y sólo si
\[
 E_M(\{m\})P=P,
 \qquad\text{equivalentemente}\qquad (M-mI)P=0.
\]
La igualdad \(PMP=mP\), por sí sola, determina una compresión escalar,
pero no garantiza que el soporte espectral preparado sea puntual.
\end{proposition}
\begin{proof}
Por el teorema espectral, \(\ker(M-mI)\) es la imagen de
\(E_M(\{m\})\). Las dos primeras condiciones expresan que
\(\operatorname{im}P\) está contenida en ese autoespacio. Para comprobar
la última distinción, tómese \(M=\operatorname{diag}(1,3)\) y el proyector
sobre \((1,1)/\sqrt2\). Entonces \(PMP=2P\), mientras que la distribución
de masa preparada tiene soporte \(\{1,3\}\), no \(\{2\}\).
\end{proof}

La individualización conserva así el alcance de cada resultado. La fibra
y el operador son objetos determinados; el paso a una línea utiliza el
proyector y la condición espectral indicados. Esta formulación se aplicará
por igual a las colecciones cargadas, neutras y compuestas.

\subsection{Leptones cargados: centro electrónico y multisecciones de generación}
\label{vi:subsec:leptones-cargados}

Las tres multisecciones cargadas poseen rangos \(1,8,16\) en profundidades
\(G_0,G_2,G_4\). Comparten la representación de carga y la trivialidad
de color, pero no el rango del soporte ni el registro de la trayectoria.
El electrón ocupa el centro \((5,5)\); las realizaciones muónica y
tauónica conservan sus multisecciones hasta efectuar el acoplamiento
espectral. La conjugación de sus representaciones se trata mediante los
operadores de la sección~\ref{vi:sec:conjugacion}, no mediante una nueva
elección de masas.

La coordenada electrónica basal \(\mathfrak e_0\) suministra la coordenatización
angular intermedia
\begin{equation}
 \mathfrak m_X^{(0)}=\mathfrak e_0
 \exp\left(n_A(X)x+\frac{s_C(X)}6y\right),
 \qquad x=A_{\mathrm{rad}},\quad y=C^*_{\mathrm{rad}}.
 \label{vi:eq:carta-angular-familias}
\end{equation}
En esta fórmula \(\mathfrak e_0\) es la sección basal y no el cierre
electrónico beta. El retorno bidireccional pertenece a la etapa posterior
de la construcción. Las dos etapas conservan su propia normalización.

Los registros nominales recuperados para las dos configuraciones no centrales
expresan
\[
 (n_A,s_C)_\mu=(42,-3),\qquad (n_A,s_C)_\tau=(64,0),
\]
y sus firmas completas declaradas son
\begin{equation}
 \sigma_\mu=(42,-3,8,0,2),\qquad
 \sigma_\tau=(64,0,25,-1,6).
 \label{vi:eq:firmas-muon-tau}
\end{equation}
Para explicar el efecto de cada coordenada, denotemos por \(\mathcal E_j\)
las lecturas sucesivas de una firma declarada:
\begin{align}
 \mathcal E_0(\sigma)&=n_Ax+s_Cy/6,\notag\\
 \mathcal E_1(\sigma)&=\mathcal E_0(\sigma)+k_{\Delta_4}\Delta_4,\notag\\
 \mathcal E_2(\sigma)&=\mathcal E_1(\sigma)+b_{120}s_{120},\notag\\
 \mathcal E_3(\sigma)&=\mathcal E_2(\sigma)+b_\zeta\zeta_{270},
 \qquad\zeta_{270}=135\Delta_4^2.
 \label{vi:eq:etapas-e3-sectoriales}
\end{align}
El símbolo \(\mathcal E_3\) designa aquí un exponente de evaluación,
no una especie ni un modo adicional. El suplemento beta que conserva esas
firmas emplea la normalización explícita
\(m_X^{[3]}=m_e^\beta\exp\mathcal E_3(\sigma_X)\).
Esta es una evaluación relativa a las firmas del suplemento; no reemplaza
la construcción ruta por ruta de \(M_X^{\beta,D}\) y
\(M_X^{\beta,\Omega}\).

La diferencia de firmas da, sin aproximación,
\begin{equation}
 \log\frac{m_\tau^{[3]}}{m_\mu^{[3]}}
 =22x+\frac12y+17\Delta_4-s_{120}+4\zeta_{270}.
 \label{vi:eq:razon-tau-muon}
\end{equation}
La demostración consiste en restar los exponentes de
\eqref{vi:eq:etapas-e3-sectoriales}. El término angular, la memoria
torsional y las dos lecturas de retorno permanecen identificados en el
resultado. Las evaluaciones conservadas son
\begin{align*}
 m_\mu^{[3]}&=105.658360294435829303\ldots\ \mathrm{MeV}/c^2,\\
 m_\tau^{[3]}&=1776.860917631432295595\ldots\ \mathrm{MeV}/c^2.
\end{align*}
Su estatuto documental es evaluación exacta al fijar la firma recuperada.
La selección física de esa firma dentro de la multisección y el criterio
de línea espectral son operaciones distintas y se conservan como tales.

En el centro electrónico el rango es uno y los lectores coinciden en
\((80,54,6)\). La historia activa
\((2,2,5,-4,-3,-2)^2\) conserva más información que su envolvente
histórica \((4,3,5,-4,-3,-5)^2\), aunque ambas compartan el patrón de
signos. Esa distinción enlaza la particularidad electrónica con la
memoria de posición de los eventos: las colecciones no se definen sólo por
la forma visual de una palabra de signos.

\subsection{Neutrinos: neutralidad, espectro de masa y transporte de sabor}
\label{vi:subsec:neutrinos-familias}

El soporte neutrínico posee dimensión \(20=2+4+6+8\), repartida entre
cuatro profundidades. Sus tres etiquetas de sabor
\(\nu_e,\nu_\mu,\nu_\tau\) y sus tres posibles autovalores de masa
pertenecen a la realización física del sector, no al censo de
profundidades. La neutralidad eléctrica conserva la orientación, la
memoria de hoja y la holonomía espinorial. Por ello no elimina las
variables que intervienen en la mezcla.

El registro neutral separa las contribuciones de dos semiciclos. Si
\(\mathcal I_m\) es el bloque de \(9\) iteraciones de la fase \(m\),
\(\theta_t\) marca la visita a \(9\), \(\epsilon_t\) conserva la hoja
y \(\omega_t\) su orientación de corona, el lector local es
\[
 T_\nu(m)=\sum_{t\in\mathcal I_m}\omega_t\epsilon_t\theta_t,
 \qquad e_m^\nu=T_\nu(m)-T_\nu(m+6).
\]
\subsubsection{Cancelaciones exactas del extractor antipodal}
\label{vi:subsubsec:cancelaciones-antipodales}

Los índices de esta subsección se leen en \(\mathbb Z/12\mathbb Z\) y
se adopta \(\operatorname{sgn}(0)=0\). Sea \(A_r\) el desplazamiento de
fase \((A_rT)(m)=T(m+r)\). Para el registro neutral,
\begin{equation}
 e^\nu=(I-A_6)T_\nu,
 \qquad e^\nu_{m+6}=-e^\nu_m.
 \label{vi:eq:extractor-antipodal}
\end{equation}
Definimos
\[
 B_a^\nu=e_a^\nu+e_{a+4}^\nu+e_{a+8}^\nu,
 \qquad
 \tau_m^\nu=\operatorname{sgn}(e_m^\nu),
 \qquad
 c_m^\nu=\tau_m^\nu\tau_{m+3}^\nu.
\]

\begin{proposition}[Anulación de los dos lectores escalares antipodales]
\label{vi:prop:cancelaciones-antipodales}
El extractor neutral satisface
\begin{equation}
 B_{a+2}^\nu=-B_a^\nu,\qquad
 b_{120}^\nu:=\sum_{a=0}^{3}\operatorname{sgn}(B_a^\nu)=0,
 \label{vi:eq:cancelacion-b120}
\end{equation}
y
\begin{equation}
 c_{m+3}^\nu=-c_m^\nu,\qquad
 b_\zeta^\nu:=\sum_{m=0}^{11}c_m^\nu=0.
 \label{vi:eq:cancelacion-bzeta}
\end{equation}
\end{proposition}

\begin{proof}
Por~\eqref{vi:eq:extractor-antipodal},
\[
 B_{a+2}^\nu
 =e_{a+2}^\nu+e_{a+6}^\nu+e_{a+10}^\nu
 =-e_{a+8}^\nu-e_a^\nu-e_{a+4}^\nu=-B_a^\nu.
\]
Los cuatro sumandos de \(b_{120}^\nu\) se cancelan en los pares
\((0,2)\) y \((1,3)\), incluso cuando alguno es nulo. Asimismo,
\(\tau_{m+6}^\nu=-\tau_m^\nu\), de donde
\[
 c_{m+3}^\nu
 =\tau_{m+3}^\nu\tau_{m+6}^\nu
 =-\tau_{m+3}^\nu\tau_m^\nu=-c_m^\nu.
\]
La suma de los doce términos se anula por pares alternos.
\end{proof}

Las cancelaciones escalares no destruyen la información orientada. En el
subespacio antiinvariante
\[
 S_-:=\ker(A_6+I),
 \qquad J:=A_3\big|_{S_-},
\]
se cumple \(J^2=-I\). Como el desplazamiento es unitario,
\(J^*=J^{-1}=-J\): el registro conserva una estructura compleja orientada.
Además,
\[
 e^\nu=(I-A_6)T_\nu=2P_-T_\nu,
 \qquad P_-=(I-A_6)/2,
\]
de modo que el vector antipodal completo permanece disponible aunque sus
dos contracciones escalares se anulen.

El retorno cíclico \(A_{12}=I\) concierne a esta representación de fase;
no equivale al retorno del estado TPK completo, que conserva acarreo, hoja,
frontera y memoria. Por la misma razón,
\eqref{vi:eq:cancelacion-b120}--\eqref{vi:eq:cancelacion-bzeta} no anulan
los coeficientes generales \(b_{120}\) y \(b_\zeta\) de una ruta no
antipodal, ni modifican las masas, las firmas compuestas o el transporte
PMNS. La costura longitudinal utilizada al prolongar el registro se
controla independientemente por el
lema~\ref{vi:lem:costura-108-1080}.


Esta diferencia mantiene la información orientada aun cuando su suma
total sea nula. El operador completo utiliza, además, los predecesores
torsionales, las vacancias y el transporte entre hojas. El paso al
subespacio físico requiere el proyector neutral de rango \(3\),
\[
 M_\nu=P_\nu\mathcal M_\nu P_\nu
 \quad\text{sobre }\operatorname{im}P_\nu.
\]
Los autovalores de esa compresión son las masas en la escala interna
acción--reloj--masa. Los ángulos de mezcla actúan entre sus marcos;
no sustituyen la escala ni generan retrospectivamente los autovalores.
La compresión coincide con la restricción del operador completo cuando
\(\operatorname{im}P_\nu\) es un subespacio reductor de
\(\mathcal M_\nu\).

El registro permite reunir un núcleo de Gram antes de diagonalizar.
Para cada tipo de dato
\(r\in\{\mathrm{car},\mathrm{hoj},\mathrm{ori},\mathrm{vac},\mathrm{hol}\}\),
sea \(\Xi_r(c)\) el vector producido por la celda o ruta \(c\).
Definimos
\begin{equation}
 K_{\mathrm{reg}}(c,d)=
 \sum_{r:N_r>0}\frac{\langle\Xi_r(c),\Xi_r(d)\rangle}{N_r},
 \qquad N_r=\sum_a\|\Xi_r(a)\|^2.
 \label{vi:eq:gram-neutrinos}
\end{equation}
La omisión de un sumando de norma total cero está fijada antes de la
evaluación. Conserva un dominio bien definido y no introduce un ajuste.

\begin{proposition}[Positividad y rango del registro neutral]
\label{vi:prop:gram-neutrinos}
El núcleo \(K_{\mathrm{reg}}\) es hermítico y positivo semidefinido.
Para una agregación tridimensional \(G\), la condición \(\det G\ne0\)
certifica que los datos agregados separan tres direcciones. Para cualquier
matriz cuadrada invertible \(G\), su factor polar
\[
 U_G=G(G^*G)^{-1/2}
\]
es unitario. Este factor no determina por sí solo los marcos espectrales
de los operadores cargado y neutral.
\end{proposition}
\begin{proof}
Para coeficientes \(z_c\), la forma cuadrática de
\eqref{vi:eq:gram-neutrinos} es
\[
 \sum_{c,d}\overline z_cK_{\mathrm{reg}}(c,d)z_d
 =\sum_{r:N_r>0}N_r^{-1}
 \left\|\sum_c z_c\Xi_r(c)\right\|^2\ge0.
\]
El rango es la dimensión de la imagen de los vectores agregados. Para la
última afirmación,
\(U_G^*U_G=(G^*G)^{-1/2}G^*G(G^*G)^{-1/2}=I\);
la invertibilidad y la igualdad de dimensiones dan también \(U_GU_G^*=I\).
La fórmula usa la matriz \(G\), mientras que un marco espectral exige
diagonalizar el operador de masa correspondiente.
\end{proof}

Si \(F_\ell,F_\nu:\mathbb C^3\to\mathcal H\) son marcos ortonormales
completos del mismo subespacio físico, entonces
\(U_{\mathrm{PMNS}}=F_\ell^*F_\nu\) es unitario. Si sus imágenes son
distintas, el producto sólo es una contracción; la unitariedad exige la
igualdad de los subespacios transportados. En presencia de degeneraciones,
la mezcla queda determinada como clase doble por los cambios de marco en
los bloques degenerados, no como una matriz numérica elegida por convenio.

La construcción mínima de fase
\[
 \Phi_K=I+(e^{i\pi/6}-1)|b_K\rangle\langle b_K|,
 \qquad\|b_K\|=1,
\]
permanece como control negativo documentado: es unitaria, pero la mejor
permutación de su identificación angular da
\(\theta_{13}\simeq17.02^\circ\), incompatible con el contraste de ese
modelo reducido. El resultado se refiere a la fase de rango uno. No
sustituye el núcleo completo \eqref{vi:eq:gram-neutrinos}, ni permite
presentar como validada una matriz de mezcla que no ha pasado su propio
control de rango y de ángulos.

\subsubsection{Autoconjugación neutral y posibilidad de un sector estéril}

La alternativa Dirac--Majorana utiliza la representación fermiónica. Una
estructura real neutral debe satisfacer
\[
 \mathcal C_\nu^2=I,\qquad
 \mathcal C_\nu M_\nu=M_\nu\mathcal C_\nu,\qquad
 \mathcal C_\nu H_\nu=H_\nu\mathcal C_\nu,
\]
con \(\mathcal C_\nu\) antiunitaria y una realificación compatible.
La conmutación con la masa preserva su parte real. Para el Hamiltoniano
autoadjunto, en cambio, la antiunitariedad conjuga también el factor
\(i\): si \(U_\nu(t)=\exp(-itH_\nu/\hbar)\), entonces
\[
 \mathcal C_\nu U_\nu(t)\mathcal C_\nu^{-1}=U_\nu(-t).
\]
Por tanto, los modos fijos \(\psi=\mathcal C_\nu\psi\) definen una
estructura real neutral, pero las condiciones escritas no bastan para
restringir a ella toda la evolución de Schr\"odinger. La conservación de
esa parte fija requiere que \(\mathcal C_\nu\) conmute con el grupo
\(U_\nu(t)\); la proposición~\ref{vi:prop:conjugacion-evolucion}
demuestra la condición correspondiente sobre su generador. Una
realización de Majorana debe declarar además la compatibilidad de la
autoconjugación con la representación fermiónica y su dinámica.
Una carga conservada no nula que cambia
de signo por \(\mathcal C_\nu\) mantiene separados los modos conjugados
en la realización de Dirac. La involución de palabra de la
sección~\ref{vi:sec:conjugacion} y la neutralidad eléctrica, consideradas
aisladamente, no deciden esa alternativa.

Un sector estéril requiere un proyector no nulo \(P_s\) con
\(P_sP_{\mathrm{act}}=0\) y \(P_sJ_{\mathrm{débil}}=0\).
Si \([P_s,M_\nu]=0\), es reductor y no mezcla con el sector activo.
Si \(P_sM_\nu P_{\mathrm{act}}\ne0\), existe mezcla y el proyector
deja de conmutar con la masa. Por tanto, una profundidad adicional del
atlas no cuenta automáticamente como cuarta especie neutrínica.

\subsection{Quarks: sabor, órbita cromática y transporte de escala}
\label{vi:subsec:quarks-familias}

El estado quark conserva la tupla
\[
 (\sigma_{ud},g,[\gamma]_{\mathrm{col}},\mathfrak M,\mathscr L),
\]
que distingue tipo arriba o abajo, generación, órbita de color,
multisección y registro. La sección~\ref{vi:sec:atlas} construye el
residuo cromático \(\kappa_{\mathrm{col}}(r,c)=r-c\pmod3\) y las
tres clases de \(12\) celdas. El índice de color actúa dentro del
soporte; el sabor conserva información de generación y firma que ese
residuo no identifica.

Los seis pares angulares nominales son
\[
\begin{array}{c|r r|r r}
 f&n_A&s_C&W_+&W_-\\\hline
 u&11&7&73&59\\
 d&17&8&110&94\\
 c&61&8&374&358\\
 s&41&-3&243&249\\
 t&100&-1&599&601\\
 b&71&-5&421&431
\end{array}
\qquad W_\pm=6n_A\pm s_C.
\]
Estos últimos enteros reescriben el mismo exponente:
\[
 n_Ax+s_Cy/6
 =\frac{W_+(x+y)+W_-(x-y)}{12}.
\]
Sumar y restar \(W_\pm\) recupera \(n_A\) y \(s_C\), de modo que
esta presentación conserva ambos canales. Las seis parejas pertenecen
a la coordenatización angular documentada; no se rellenan con ceros las tres
coordenadas restantes para fabricar una firma completa. Sus operadores
de color y las evaluaciones angulares se preservan como resultados de
distinto alcance.

\begin{proposition}[Degeneración cromática de una masa de sabor]
\label{vi:prop:masa-color}
Sea \(V_{\mathrm{col}}\simeq\mathbb C^3\) la representación cromática
fundamental y \(U\) su acción unitaria irreducible. La reducción
\[
 \mathbb E_{\mathrm{col}}(M)=
 \int_{SU(3)}U(g)MU(g)^*\,dg
\]
es un múltiplo de \(I_3\). Si \(M\) ya conmuta con la acción, sus tres
componentes cromáticas tienen el mismo valor de masa de sabor.
\end{proposition}
\begin{proof}
La invariancia de la medida de Haar da
\(U(h)\mathbb E_{\mathrm{col}}(M)U(h)^*=
\mathbb E_{\mathrm{col}}(M)\). El lema de Schur hace escalar este
operador. Como la traza se conserva, el escalar es
\(\operatorname{Tr}M/3\). Si \(M\) es invariante desde el inicio,
la integral devuelve \(M\).
\end{proof}

La prueba explica qué elimina el cociente de color: la dependencia de
la cifra respecto del representante rojo, verde o azul. No elimina las
rutas cromáticas que intervienen después en un estado ligado. Tampoco
convierte esa reducción en un selector de sabor.

La masa corriente se presenta como sección en una escala y un esquema
declarados, \(m_q(\mu,\mathscr R)\). Si \(t=\log\mu\) y la conexión
induce \(\gamma_m(g(t))\), el transporte satisface
\[
 \left(\frac{d}{dt}+\gamma_m(g(t))\right)m_q=0,
 \qquad
 m_q(t_2)=\mathcal P\exp\left(-\int_{t_1}^{t_2}
 \gamma_m(g(t))\,dt\right)m_q(t_1).
\]
La conexión y su endomorfismo son datos del transporte, no números
elegidos para igualar una tabla. Un residual físico sólo tiene sentido
después de situar ambas secciones en la misma escala y el mismo esquema.

La mezcla CKM actúa sobre otra coordenada: el marco de sabor. Una vez
construida la matriz unitaria interna, el operador abajo, expresado en
el marco arriba, es \(B_d^{(u)}=V_{\mathrm{CKM}}B_dV_{\mathrm{CKM}}^*\).
La conjugación unitaria conserva el polinomio característico, pues
\[
 \det(zI-VB_dV^*)=\det\bigl(V(zI-B_d)V^*\bigr)=\det(zI-B_d).
\]
Así se entiende la función de los ángulos: relacionan marcos, sin
convertirse en generadores retrospectivos de los autovalores de masa.

\subsection{Coordenatizaciones de masa nula: fotón y representación adjunta de color}
\label{vi:subsec:nulo-familias}

El dominio nulo se separa antes de aplicar el carácter positivo. En la
realización total,
\[
 \mathcal H_{\mathrm{phys}}=\mathcal H_0\oplus\mathcal H_+,
 \qquad M|_{\mathcal H_0}=0,\qquad M|_{\mathcal H_+}>0.
\]
La exponencial de una firma finita no se anula. Por eso una coordenatización nula
no se obtiene asignando a una ruta positiva una firma arbitrariamente
grande y llamando cero a su aproximación.

El fotón pertenece a la coordenatización electromagnética, con su representación
de calibre y sus dos polarizaciones transversales físicas. Los gluones
pertenecen a la representación adjunta cromática, de dimensión \(8\).
La igualdad de sus masas nulas no identifica ambas representaciones.
La fibra, el calibre, las polarizaciones y los canales de acoplamiento
siguen formando parte del estado.

También es compatible una propagación orientada con masa nula. En la
coordenatización documentada,
\[
 c_\pm=\frac{c}{1\pm\varepsilon},\qquad
 \varepsilon=q_-^{90}-q_+^{90},\qquad0<\varepsilon<1.
\]
Entonces \(c_\pm>0\) y
\[
 \frac{2}{c_+^{-1}+c_-^{-1}}=c.
\]
La identidad procede de sumar \((1+\varepsilon)/c\) y
\((1-\varepsilon)/c\). La asimetría de las dos lecturas conserva su
información de orientación; no expresa por sí misma un término de masa.

\subsection{Bosones masivos y sector escalar}
\label{vi:subsec:bosones-familias}

Los descriptores de \(W\), \(Z\) y del sector escalar no se fuerzan
dentro de una de las \(13\) multisecciones fermiónicas. Conservan su
representación y su dominio de realización. Las lecturas nominales
documentadas son
\[
\begin{array}{c|c|c|c|c}
 X&\text{ruta nominal}&\text{dirección}&(n_A,s_C)&(W_+,W_-)\\\hline
 W&15/37&E\text{--}O&(94,-1)&(563,565)\\
 Z&20/23&E\text{--}O&(95,-1)&(569,571)\\
 H&24/29&N\text{--}S&(97,9)&(591,573)
\end{array}
\]
Estos nombres de ruta pertenecen al registro nominal conservado. Su
unión con el dominio orientado completo debe retener los datos de
trayectoria, representación y proyector. Los recuentos cardinales
expresados, respectivamente \((23,44,11,30)\), \((33,45,10,20)\) y
\((40,34,22,12)\), se conservan junto a la dirección y no se reemplazan
por una sola etiqueta.

Para \(W\) y \(Z\) se recuperan las firmas completas
\[
 \sigma_W=(94,-1,0,-2,4),\qquad
 \sigma_Z=(95,-1,-11,0,-4).
\]
La lectura \eqref{vi:eq:etapas-e3-sectoriales} produce, para esas firmas,
\[
 m_W^{[3]}=80378.964909704547024865\ldots\ \mathrm{MeV}/c^2,
\]
\[
 m_Z^{[3]}=91187.533444313407844855\ldots\ \mathrm{MeV}/c^2.
\]
Su razón conserva exactamente la diferencia de registros evaluados:
\[
 \log\frac{m_Z^{[3]}}{m_W^{[3]}}
 =x-11\Delta_4+2s_{120}-8\zeta_{270}.
\]
En cambio, la coordenada nominal \((97,9)\) del sector escalar pertenece
a la coordenatización angular \(\mathbb Z^2\) documentada: no se le atribuye por
esta escritura una firma \(\mathbb Z^5\) ni un refinamiento de orden
\(3\). Su realización como excitación escalar exige el proyector
escalar y los acoplamientos de calibre y fermiónicos correspondientes.

Para una especie inestable, el resultado físico es un polo del operador
de canales, no únicamente la evaluación de un carácter positivo. La
subsección~\ref{vi:subsec:resonancias-familias} construye esa lectura y
separa masa de polo y anchura. En el registro documental conservado,
las cifras anteriores son evaluaciones condicionadas a la firma; la
individualización de ruta, proyector y polo constituye la realización
sectorial que debe acompañarlas.

\subsection{Gramática finita de las composiciones materiales}
\label{vi:subsec:gramatica-compuestos}

Sean \(\mathcal Q_a\), \(a\in\mathbb Z/3\mathbb Z\), las tres
clases de color del atlas, cada una de cardinal \(12\). El dual formal
\(q^\vee\) conserva el soporte conjugado y cambia la orientación de
carga y color. No es la mera celda \(\rho_{\mathrm c}(q)\) desprovista
de esos datos. Las expresiones mesónicas y bariónicas se definen por
\begin{align}
 \mathfrak M&=\{q\otimes(q')^\vee:
 q,q'\in\mathcal Q_a\text{ para algún }a\},\notag\\
 \mathfrak B&=\mathcal Q_0\times\mathcal Q_1\times\mathcal Q_2.
 \label{vi:eq:gramatica-meson-barion}
\end{align}
Los pares son ordenados y las posiciones de color del triple están
fijadas. La neutralidad residual vale, respectivamente,
\(a-a=0\) y \(0+1+2=0\pmod3\).

\begin{proposition}[Censo composicional interno]
\label{vi:prop:censo-compuestos}
La gramática contiene \(432\) expresiones mesónicas, \(1728\)
expresiones bariónicas y \(372\) flechas cargadas compatibles, repartidas
en \(320\) leptónicas y \(52\) quark. Estos cardinales cuentan
expresiones y flechas internas, no especies físicas ni predicciones
independientes de masas.
\end{proposition}
\begin{proof}
Para cada color se elige un par ordenado de \(12\) celdas; por tanto,
\(|\mathfrak M|=3\cdot12^2=432\). Las tres posiciones del barión dan
\(|\mathfrak B|=12^3=1728\). Las transiciones leptón--neutrino de
igual profundidad sólo comparten \(G_2,G_4\); en ambas direcciones se
obtiene \(2(8\cdot4+16\cdot8)=320\).

Para las transiciones arriba--abajo se exige además igualdad de color.
En \(G_1\), las multiplicidades cromáticas son \((2,1,1)\) y
\((0,1,1)\); en \(G_3\), son \((4,4,4)\) y \((2,2,2)\).
Hay \(2+24=26\) pares de una dirección y \(52\) de ambas. Sumar
\(320+52\) completa el censo. Todos los factores proceden del atlas
anterior a la evaluación de masas.
\end{proof}

Las \(81\) identidades de las celdas completan las flechas neutras
de esta gramática, pero una identidad composicional no se convierte por
ese hecho en un bosón neutro. La representación y el canal de realización
son los datos que distinguen ambas nociones.

\subsection{Composición del registro y memoria de enlace}
\label{vi:subsec:composicion-registro}

Un compuesto contiene más información que la lista de sus
constituyentes. Sean \(L_1,L_2\) registros enriquecidos, \(g\) el
transporte que alinea fase, hoja y orientación, y \(B\) la subtraza
identificada en la frontera común. En las coordenadas de acción y
eventos, la composición es
\begin{equation}
 (n,e)_1\star_{(g,B)}(n,e)_2
 =\bigl(n_1+n_2-n_B,\ e_1+ge_2-e_B\bigr).
 \label{vi:eq:composicion-frontera}
\end{equation}
La sustracción evita contar dos veces una misma subtraza; el transporte
conserva la orientación en la interfaz. Los predecesores torsionales y
las marcas de corona se componen en el mismo nivel antes de volver a
aplicar \(L_{\mathrm{tick}}\). La descomposición reversible de eventos
de la proposición~\ref{vi:prop:eventos} permite transportar el registro
completo sin reducirlo a su total.

Para un árbol ordenado \(\mathcal T\) de \(n\) constituyentes se obtiene
\begin{equation}
 (\Gamma_1,\ldots,\Gamma_n;\mathcal T,\partial)
 \xmapsto{\mathfrak C_{12}^{(n)}}L_Y
 \xmapsto{L_{\mathrm{tick}}}\sigma_Y
 \longmapsto (\eta_Y,M_Y).
 \label{vi:eq:cadena-compuesto}
\end{equation}
La frontera \(\partial\) incluye los transportes y subtrazas de cada
arista del árbol. Una vez fijados esos datos admisibles, cada composición
es una operación explícita. El carácter se aplica a la firma final;
no selecciona previamente el árbol por la masa de llegada.

\begin{proposition}[Compatibilidad de la parentización]
\label{vi:prop:cociclo-compuesto}
En una coordenatización de registros donde
\(L_1\star L_2=L_1+L_2+\mathfrak b(L_1,L_2)\), las dos
parentizaciones de un triple admisible coinciden si y sólo si
\begin{align*}
 \mathfrak b(L_1,L_2)+\mathfrak b(L_1\star L_2,L_3)
 =\mathfrak b(L_2,L_3)+\mathfrak b(L_1,L_2\star L_3).
\end{align*}
La igualdad se exige en el dominio donde ambos árboles representan la
misma interfaz orientada.
\end{proposition}
\begin{proof}
Desarrollar \((L_1\star L_2)\star L_3\) da
\(L_1+L_2+L_3+\mathfrak b(L_1,L_2)+\mathfrak b(L_1\star L_2,L_3)\).
Desarrollar \(L_1\star(L_2\star L_3)\) da la misma suma de tres
registros y los dos términos del segundo miembro. Su igualdad equivale
exactamente a la identidad enunciada.
\end{proof}

La ley cocíclica asegura independencia de una parentización equivalente;
no identifica árboles físicamente distintos. Dos estructuras de enlace
que conservan canales diferentes permanecen como realizaciones distintas
hasta disponer de un entrelazador de registros, proyectores y operadores.

\subsubsection{Un testigo finito de la información que perdería una suma de firmas}

El extractor reducido histórico de eventos se define por
\[
 F_0(n,e)=\left(n,\sum_m e_m,K_0(e),
 \sum_{a=0}^3\operatorname{sgn}(e_a+e_{a+4}+e_{a+8}),
 \sum_m\tau_m\tau_{m+3}\right),
\]
donde \(\tau_m=\operatorname{sgn}(e_m)\) y
\(K_0(e)=\sum_{\tau_m=0}\Sigma_{12}(m)\), con
\(\Sigma_{12}=(+,+,+,-,-,-)^2\). Su tercera coordenada es una
lectura de eventos nulos; no es la torsión fuerte
\(k_{\Delta_4}=T_z-2C_z\) de la
sección~\ref{vi:sec:operador-masa}.

Considérense los registros de las rutas
\[
 a=(1,3,-1,+1),\quad b=(6,7,-1,+1),\quad c=(1,2,-1,-1),
\]
con
\begin{align*}
 e(a)&=(5,4,4,-4,-4,-6)^2,\\
 e(b)&=(4,5,4,-4,-6,-4)^2,\\
 e(c)&=(1,4,1,-2,-4,-3)^2.
\end{align*}
Tienen \(F_0(a)=F_0(b)=(8,-2,0,0,-12)\) y
\(F_0(c)=(28,-6,0,-4,-12)\). Al superponer los registros,
\[
 F_0(a\boxplus c)=(36,-8,0,0,-12),\qquad
 F_0(b\boxplus c)=(36,-8,0,-2,-12).
\]
Las entradas proyectadas son iguales y las salidas compuestas difieren.
Por consiguiente, ninguna operación sobre las dos cinco-tuplas de esta
proyección recupera todas las superposiciones. El ejemplo demuestra una
obstrucción precisa para \(F_0\); no atribuye el mismo resultado a toda
posible ampliación de una firma. La construcción vigente evita esa
pérdida mediante \eqref{vi:eq:cadena-compuesto}: conserva los eventos
y los predecesores, compone y sólo entonces evalúa el extractor completo.

\subsection{Mesones: pareja dual, singlete y diferenciación de las excitaciones}
\label{vi:subsec:mesones-familias}

La expresión \(q\otimes(q')^\vee\) establece el par ordenado y su
neutralidad residual. El singlete de color pertenece a la realización
\(V\otimes\overline V\), \(V=\mathbb C^3\). Si
\(\Omega=\sum_{a=1}^3|a\bar a\rangle\), el proyector es
\begin{equation}
 P_M^{(1)}=\frac13|\Omega\rangle\langle\Omega|.
 \label{vi:eq:singlete-meson}
\end{equation}

\begin{proposition}[Proyección mesónica de color]
La aplicación \(P_M^{(1)}\) es un proyector ortogonal de rango uno y
su imagen es invariante bajo \(U\otimes\overline U\), \(U\in SU(3)\).
\end{proposition}
\begin{proof}
Se tiene \(\langle\Omega,\Omega\rangle=3\), luego
\((P_M^{(1)})^2=P_M^{(1)}\). Además,
\[
 (U\otimes\overline U)\Omega
 =\sum_{a,b,c}U_{ba}\overline U_{ca}|b\bar c\rangle
 =\sum_b|b\bar b\rangle=\Omega.
\]
La imagen conserva así la contracción cromática sin privilegiar un color.
\end{proof}

El descriptor mesónico reúne sabores, espín y paridad, isospín y carga,
excitación radial u orbital, árbol de enlace y frontera. La etiqueta
\(J^{PC}\) se utiliza cuando \(C\) está definido para el canal; no se
asigna automáticamente conjugación de carga a un estado cargado. Dos
estados con los mismos sabores pueden tener distintos proyectores
radiales, orbitales o de canal y, por tanto, espectros distintos. La
construcción orbital de la sección~\ref{vi:sec:kepler} aporta una
lectura posible cuando se cumplen sus hipótesis dinámicas; los números
orbitales no se deducen del solo nombre del mesón.

El registro de evaluaciones compuestas conserva
\[
\begin{array}{c|r r r r r}
 &n_A&s_C&k_{\Delta_4}&b_{120}&b_\zeta\\\hline
 \pi^0&44&-4&-25&1&-2\\
 \pi^\pm&44&1&-2&2&-1\\
 K^\pm&54&-1&17&-2&2\\
 K^0&54&0&32&3&-2
\end{array}
\]
Para estas firmas, la misma lectura \eqref{vi:eq:etapas-e3-sectoriales}
explica las diferencias de los pares:
\begin{align}
 \log\frac{m_{\pi^\pm}^{[3]}}{m_{\pi^0}^{[3]}}
 &=\frac56y+23\Delta_4+s_{120}+\zeta_{270},\notag\\
 \log\frac{m_{K^0}^{[3]}}{m_{K^\pm}^{[3]}}
 &=\frac16y+15\Delta_4+5s_{120}-4\zeta_{270}.
 \label{vi:eq:desdoblamientos-mesonicos}
\end{align}
Los términos de acción \(44x\) y \(54x\) se cancelan dentro de
cada pareja. Las diferencias restantes conservan qué coordenada de
registro produce cada contribución. Son consecuencias exactas de las
firmas declaradas, no una prueba sustituta de su generación por el árbol
de enlace.

Las evaluaciones resultantes son
\begin{align*}
 m_{\pi^0}^{[3]}&=134.976724327872125739\ldots\ \mathrm{MeV}/c^2,\\
 m_{\pi^\pm}^{[3]}&=139.570485171572191374\ldots\ \mathrm{MeV}/c^2,\\
 m_{K^\pm}^{[3]}&=493.677085649530612475\ldots\ \mathrm{MeV}/c^2,\\
 m_{K^0}^{[3]}&=497.611186497750457358\ldots\ \mathrm{MeV}/c^2.
\end{align*}
Cada cifra conserva el
estatuto de evaluación exacta condicionada a la firma compuesta
recuperada. El singlete, la gramática y la extracción del registro están
construidos; la realización nominal completa exige que el árbol de
ligadura regenere esa firma antes de abrir el contraste.

\subsection{Bariones: triple cromático, intercambio y diferencia neutrón--protón}
\label{vi:subsec:bariones-familias}

La neutralidad de \(q_0\otimes q_1\otimes q_2\) conserva las tres
posiciones cromáticas. En la realización tensorial se introduce el vector
normalizado
\[
 |\varepsilon\rangle=\frac1{\sqrt6}
 \sum_{a,b,c=1}^3\varepsilon_{abc}|abc\rangle,
 \qquad P_B^{(1)}=|\varepsilon\rangle\langle\varepsilon|.
\]

\begin{proposition}[Singlete bariónico y regla de intercambio]
\label{vi:prop:singlete-barion}
El operador \(P_B^{(1)}\) es un proyector ortogonal de rango uno,
invariante bajo \(U^{\otimes3}\), \(U\in SU(3)\). Una permutación
de dos posiciones cambia el signo de su vector generador y conserva
el proyector.
\end{proposition}
\begin{proof}
Hay \(6\) términos no nulos ortogonales, cada uno de módulo
\(1/\sqrt6\); por tanto el vector tiene norma uno. La identidad
del determinante da
\(U^{\otimes3}|\varepsilon\rangle=
\det(U)|\varepsilon\rangle=|\varepsilon\rangle\).
Una transposición cambia el signo de \(\varepsilon_{abc}\). En el
producto exterior del vector por su adjunto, ambos signos se cancelan.
\end{proof}

La antisimetría cromática es una parte del estado completo. El proyector
físico reúne color, espín y paridad, isospín, carga y simetría de
intercambio de las componentes orbitales y de sabor. El principio de
Pauli actúa en el producto fermiónico completo; no basta aplicarlo a una
lista desnuda de colores. El árbol triple registra qué par se acopló
primero y cómo atravesó la frontera el tercer constituyente. La
proposición~\ref{vi:prop:cociclo-compuesto} controla las
parentizaciones equivalentes.

Las firmas recuperadas de protón y neutrón son
\[
 \sigma_p=(59,0,9,1,0),\qquad
 \sigma_n=(59,1,-34,0,1).
\]
Comparten la coordenada de acción \(59\), pero difieren en hoja,
torsión y retornos. Su diferencia de exponente es
\begin{equation}
 \log\frac{m_n^{[3]}}{m_p^{[3]}}
 =\frac16y-43\Delta_4-s_{120}+\zeta_{270}.
 \label{vi:eq:diferencia-neutron-proton}
\end{equation}
La igualdad se demuestra restando las cinco coordenadas. En esta firma,
\(-34\) es un entero torsional del neutrón; no se identifica con el
exponente decimal \(-34\) de la escala absoluta de acción. La igualdad
de dos enteros en mapas distintos no elimina su tipo ni su procedencia.

Las evaluaciones conservadas son
\begin{align*}
 m_p^{[3]}&=938.271751546984898298\ldots\ \mathrm{MeV}/c^2,\\
 m_n^{[3]}&=939.565810472507023312\ldots\ \mathrm{MeV}/c^2.
\end{align*}
La fórmula \eqref{vi:eq:diferencia-neutron-proton} localiza el efecto de
cada componente de las firmas nominales. La derivación física de la
diferencia exige transportar isospín, orientación de vacancias,
frontera de enlace y contribución electromagnética a una misma coordenatización.
La evaluación del carácter conserva su exactitud bajo la firma fijada;
esa exactitud no sustituye la construcción prospectiva del enlace.

\subsection{Composiciones de orden superior y equivalencia de canales}
\label{vi:subsec:compuestos-superiores}

La aplicación \(\mathfrak C_{12}^{(n)}\) admite un número arbitrario
finito de constituyentes. Un tetraquark emplea \(qq\bar q\bar q\);
un pentaquark, \(qqqq\bar q\); un estado gluónico ligado usa
representaciones adjuntas proyectadas al singlete; un híbrido conserva
constituyentes fermiónicos y de calibre en el mismo árbol. Núcleos,
átomos y moléculas emplean la misma precedencia de registro y lectura,
con sus propios acoplamientos y fronteras.

La energía de enlace pertenece a la realización del registro compuesto.
No se obtiene sumando masas aisladas y bautizando como ligadura el
residual deseado. Incluso cuando un cociclo toma valores aditivos sobre
registros, el extractor de signos y correlaciones se aplica después;
por ello el carácter multiplicativo de una firma no transforma la masa
compuesta en una suma de masas constituyentes.

Si dos árboles \(\mathcal T_1,\mathcal T_2\) poseen un unitario
\(V\) que entrelaza sus operadores y sus proyectores de canal,
\[
 VH_1=H_2V,\qquad VP_1=P_2V,
\]
entonces \(V(z-H_1)^{-1}=(z-H_2)^{-1}V\) en el resolvente común.
Su continuación, cuando existe y se transporta en la misma coordenatización,
conserva los polos y sus multiplicidades. Esta igualdad da el criterio
operativo para identificar dos realizaciones de un compuesto; compartir
sabores o una masa aproximada no basta.

\subsection{Resonancias: masa de polo, anchura y memoria de supervivencia}
\label{vi:subsec:resonancias-familias}

Sea \(P\) el subespacio preparado y \(Q=I-P\) el de canales
eliminados. En una región donde \(z-QHQ\) es invertible, la eliminación
de la componente \(Q\psi\) de \((z-H)\psi=0\) da
\begin{equation}
 H_{\mathrm{eff}}(z)=PHP+PHQ(z-QHQ)^{-1}QHP.
 \label{vi:eq:feshbach-familias}
\end{equation}

\begin{proposition}[Lectura reducida de un estado ligado o resonante]
\label{vi:prop:feshbach-familias}
Si \(P\mathcal H\) es finito-dimensional, los valores espectrales
admisibles del problema reducido satisfacen
\(\det(z-H_{\mathrm{eff}}(z))=0\). Una resonancia exige, además,
continuación meromorfa a la coordenatización de canales pertinente y un cero
aislado con residuo físico admisible.
\end{proposition}
\begin{proof}
La ecuación en el complemento es
\((z-QHQ)Q\psi=QHP\psi\). Resolverla y sustituir en la ecuación
de \(P\psi\) produce \((z-H_{\mathrm{eff}}(z))P\psi=0\).
En dimensión finita existe solución no nula exactamente cuando se
anula el determinante. El cálculo algebraico se realiza en el dominio
del resolvente; el paso a una resonancia utiliza la continuación y el
residuo indicados, que no se deducen de la sola apertura de un canal.
\end{proof}

Un estado estable es una línea espectral normalizable sin canal de
decaimiento abierto; un estado ligado está por debajo del umbral
pertinente, aunque pueda decaer por otro canal; una resonancia es un
polo de la continuación. Si
\[
 z_X=M_Xc_{\mathrm{int}}^2-\frac{i}{2}\Gamma_X,
 \qquad\Gamma_X>0,
\]
la parte real y la imaginaria son dos lectores coordinados del mismo
polo. La anchura no es una sexta coordenada de \(\sigma\).

La lectura discreta usa un autovalor de supervivencia
\(\lambda_X=re^{-i\theta}\), \(0<r<1\), y un paso temporal
\(t_0>0\). Fijada la rama \(\theta\in I_\theta\) de la coordenatización de
cuasienergía,
\begin{equation}
 z_X=\frac{\hbar_{\mathrm{int}}}{t_0}
 (\theta+i\log r),\quad
 M_X=\frac{\hbar_{\mathrm{int}}\theta}{t_0c_{\mathrm{int}}^2},\quad
 \Gamma_X=-\frac{2\hbar_{\mathrm{int}}}{t_0}\log r.
 \label{vi:eq:anchura-discreta-familias}
\end{equation}
En efecto, sustituir en
\(\lambda_X=\exp(-iz_Xt_0/\hbar_{\mathrm{int}})\) recupera
\(re^{-i\theta}\). Como \(\log r<0\), el polo es retardado y la
anchura positiva. La probabilidad de supervivencia tras \(N\) pasos es
\(r^{2N}=\exp(-\Gamma_XNt_0/\hbar_{\mathrm{int}})\), de donde
\(\tau_X=\hbar_{\mathrm{int}}/\Gamma_X\).

Si esta lectura se obtiene de una compresión \(SUS\), la interpretación
asintótica mediante un único modo requiere un autovalor dominante
simple, solapamiento inicial no nulo y separación en módulo del resto
del espectro. Una normalización sobre probabilidades, en lugar de
amplitudes, modifica el factor \(2\). La rama temporal y esa convención
acompañan por tanto a cada anchura expresada.

La igualdad de masa y anchura de estados conjugados requiere transportar
también los canales:
\[
 U_C H_{\mathrm{eff}}^X(z)U_C^{-1}=H_{\mathrm{eff}}^{\bar X}(z).
\]
Entonces se transportan los mismos polos con la misma prescripción
causal. Una representación antiunitaria exige transportar además esa
prescripción: la conjugación desnuda de \(z\) intercambia polos
retardados y avanzados.

\subsection{Tres codominios topológicos y sector virtual}
\label{vi:subsec:topologicos-familias}

Las clases topológicas conservan objetos y reglas de fusión, no una
masa escalar universal. El dato primario es
\((N_{ab}^{c},F^{abc}_d,R^{ab}_c)\), con sus ecuaciones de coherencia.
Una brecha de energía requiere después una realización Hamiltoniana y
un proyector espectral. La asignación de ese tipo de salida permite
incluir estos sectores sin convertirlos en filas de masa elemental.

\subsubsection{Sector Fibonacci}

Los objetos \(\mathbf1,\tau\) satisfacen
\(\tau\otimes\tau=\mathbf1\oplus\tau\). La dimensión positiva
\(d_\tau\) conserva la multiplicación de fusión:
\(d_\tau^2=1+d_\tau\). La realización del corpus identifica esa
dimensión con la coordenada interna \(\varphi\) ya generada; la
ecuación de fusión es su reconocimiento en este codominio, no un
sustituto de su genealogía APP--TRIT--TPK. En la convención fijada,
\[
 F_\tau^{\tau\tau\tau}=\begin{pmatrix}
 \varphi^{-1}&\varphi^{-1/2}\\
 \varphi^{-1/2}&-\varphi^{-1}
 \end{pmatrix},\qquad
 R_{\mathbf1}^{\tau\tau}=e^{-4\pi i/5},\quad
 R_\tau^{\tau\tau}=e^{3\pi i/5}.
\]
La identidad \(\varphi^{-2}+\varphi^{-1}=1\) prueba directamente
que \(F^*F=I\); los términos no diagonales se cancelan. Las fases
de \(R\) conservan la orientación del trenzado.

\subsubsection{Sector Ising y realización Majorana}

Las reglas son
\[
 \psi\otimes\psi=\mathbf1,\quad
 \psi\otimes\sigma=\sigma,\quad
 \sigma\otimes\sigma=\mathbf1\oplus\psi,
\]
con
\[
 F_\sigma^{\sigma\sigma\sigma}=\frac1{\sqrt2}
 \begin{pmatrix}1&1\\1&-1\end{pmatrix},\qquad
 R_{\mathbf1}^{\sigma\sigma}=e^{-\pi i/8},\quad
 R_\psi^{\sigma\sigma}=e^{3\pi i/8}.
\]
La multiplicación matricial da \(F^2=I\). El sector conserva sus
tres tipos de fusión y su representación; su denominación Majorana
no convierte automáticamente un neutrino del sector anterior en modo
autoconjugado. Esa afirmación utiliza las condiciones antiunitarias
específicas de la realización neutral.

\subsubsection{Sector abeliano de seis clases}

Para \(a,b\in\mathbb Z/6\mathbb Z\),
\[
 a\otimes b=a+b\pmod6,\qquad
 R_{ab}=\zeta_6^{ab},\qquad\zeta_6=e^{2\pi i/6}.
\]
La identidad \(R_{a+b,c}=R_{a,c}R_{b,c}\) y su análoga en el
segundo argumento expresan el carácter bilineal del trenzado.
La construcción duplicada
\(\mathcal C\boxtimes\mathcal C^{\mathrm{rev}}\) conserva ambas
orientaciones. Sus clases no deben contarse como seis nuevas masas
elementales.

\subsubsection{Persistencia finita y virtualidad}

El sector virtual conserva extensiones finitas, su canal y su horizonte
de persistencia. Una sección finita con obstrucción terminal no ocupa
el mismo codominio que una sección asintótica persistente. La
virtualidad se lee en una amplitud o resolvente y no equivale, por
definición, a una masa imaginaria universal. Tampoco basta que una
extensión global no esté certificada en una ficha para afirmar que
existe una obstrucción: el horizonte y su prueba son datos diferentes.
La clasificación conserva precisamente esa información finita sin
inventar una partícula estable para cada línea interna de cálculo.

\subsection{Clasificación de veintitrés sectores de partículas y alcance de sus realizaciones}
\label{vi:subsec:23-clases}

El inventario sectorial reúne \(17\) clases elementales, \(2\) clases
compuestas, \(3\) codominios topológicos y \(1\) sector virtual:
\(17+2+3+1=23\). Las tres filas neutrínicas nombran sabores, las
seis filas quark conservan su triplicidad cromática y la fila gluónica
conserva la representación adjunta. Ninguna de estas convenciones de
conteo modifica el atlas \(81/56/13/324\).

La siguiente relación mantiene el tipo de salida y el alcance de las
realizaciones documentadas. Las expresiones «firma declarada» y
«operador de fibra» remiten a las construcciones explícitas de esta
sección, no a grados distintos de precisión decimal.
El registro de procedencia distingue el cierre escalar electrónico y
las dos realizaciones nulas de la escalarización física todavía pendiente en
las configuraciones no centrales. En los neutrinos conserva pendiente la
realización completa de mezcla y escala; en los quarks, la selección de
sabor y su coordenatización de escala; en los bosones nominales, la unión de ruta
y polo; en los compuestos nombrados, la realización escalar de su
enlace. Estos estados documentales no alteran los operadores, las
gramáticas ni las evaluaciones condicionadas que se han expuesto.

\begin{table}[!htbp]
\centering\fontsize{9.2}{10.5}\selectfont
\setlength{\tabcolsep}{4pt}\renewcommand{\arraystretch}{1.04}
\begin{tabular}{@{}p{.11\linewidth}p{.20\linewidth}p{.62\linewidth}@{}}
 \toprule
 Clase&Representantes&Resultado conservado y operación de realización\\
 \midrule
 
 1&\(e^-,e^+\)&Cierre beta escalar de rango uno y representación conjugada.\\
 2&\(\mu^-,\mu^+\)&Operadores exactos de fibra; evaluación nominal y refinamiento de orden \(3\) con firma declarada; selección espectral física diferenciada.\\
 3&\(\tau^-,\tau^+\)&Operadores exactos de fibra; evaluación nominal y refinamiento de orden \(3\) con firma declarada; selección espectral física diferenciada.\\
 4&\(\nu_e,\bar\nu_e\)&Operador neutral de fibra; realización de mezcla, marco de masa y escala física por sus mapas completos.\\
 5&\(\nu_\mu,\bar\nu_\mu\)&Mismo dominio neutral, con etiqueta de sabor distinta; no una profundidad individual del atlas.\\
 6&\(\nu_\tau,\bar\nu_\tau\)&Mismo dominio neutral; conservación de la alternativa Dirac--Majorana bajo sus condiciones.\\
 7&\(u,\bar u\)&Operador de fibra cromática y lectura angular nominal; escalarización, esquema y escala conservados.\\
 8&\(d,\bar d\)&Operador de fibra cromática y lectura angular nominal; escalarización, esquema y escala conservados.\\
 9&\(c,\bar c\)&Operador de fibra cromática y lectura angular nominal; escalarización, esquema y escala conservados.\\
 10&\(s,\bar s\)&Operador de fibra cromática y lectura angular nominal; escalarización, esquema y escala conservados.\\
 11&\(t,\bar t\)&Operador de fibra cromática y lectura angular nominal; convención de masa y canal de comparación explícitos.\\
 12&\(b,\bar b\)&Operador de fibra cromática y lectura angular nominal; escalarización, esquema y escala conservados.\\
 13&Fotón&Sección de masa exactamente nula; representación electromagnética.\\
 14&Gluones&Sección de masa exactamente nula; representación adjunta de color.\\
 15&\(W^+,W^-\)&Firma nominal completa y evaluación de orden \(3\); unión a ruta, proyector y polo como realización física diferenciada.\\
 16&\(Z^0\)&Firma nominal completa y evaluación de orden \(3\); unión a ruta, proyector y polo como realización física diferenciada.\\
 17&\(H\)&Descriptor escalar y evaluación angular nominal; proyector, acoplamientos y polo de la realización escalar.\\
 18&Mesones&Gramática exacta de pareja dual y enlace; firmas compuestas recuperadas; escalarización nominal y canales de polo diferenciados.\\
 19&Bariones&Gramática exacta de triple cromático y enlace; firmas de protón y neutrón recuperadas; realización nominal de ligadura diferenciada.\\
 20&Fibonacci&Codominio de fusión y trenzado; una brecha requiere su realización Hamiltoniana.\\
 21&Ising/Majorana&Codominio de fusión y estructura real según la representación; sin masa escalar universal asignada.\\
 22&\(\mathbb Z/6\mathbb Z\)&Seis clases de trenzado; codominio topológico, no seis masas universales.\\
 23&Secciones virtuales&Extensión finita y horizonte de obstrucción; sin escalar estable universal.\\
\end{tabular}
\end{table}

Los estados documentales exactos se expresan en el apéndice~B.3: la
tabla~\ref{tab:vi-estados-clases} conserva la leyenda de estatutos y la
tabla~\ref{tab:vi-clases} conserva, para las veintitrés clases, residencia,
incidencia, representación y estado escalar.

\begin{proposition}[Producto fibrado de realización y codominio observable]
\label{vi:prop:producto-fibrado-especie-ruta}
Para \(X\in\mathcal P_{\rm adm}\), la fibra definida en
\eqref{vi:eq:fibra-descriptor} es no vacía si y sólo si
\[
 d_{\mathcal P}(X)\in\operatorname{im}d_{\mathcal R}.
\]
Sobre esa fibra, la medida espectral conjunta
\eqref{vi:eq:familias-pvm} determina el codominio observable de los lectores
\((M_X^{\beta,D},M_X^{\beta,\Omega})\). Una preparación \(P_X\) expresa una
línea escalar \(m\) exactamente cuando satisface el criterio de la
proposición~\ref{vi:prop:familias-linea}; en los demás casos conserva el
multiplete, la distribución espectral o el tipo sectorial producido.
\end{proposition}
\begin{proof}
Por la definición de imagen, la condición escrita equivale a la existencia de
una ruta \(\gamma\in\mathcal R_{\rm adm}\) con
\(d_{\mathcal R}(\gamma)=d_{\mathcal P}(X)\); ésta es precisamente la
condición de pertenencia a \(F_X\). La segunda afirmación sigue del teorema
espectral aplicado a los dos operadores diagonales y de la caracterización
del soporte puntual demostrada en la
proposición~\ref{vi:prop:familias-linea}.
\end{proof}

El corte documental externo contiene \(471\) observables de masa y
\(384\) de anchura, en total \(855\). Son filas de contraste, no
\(855\) predicciones internas. En particular, las \(243\) filas
mesónicas de masa y \(166\) bariónicas no coinciden con los cardinales
\(432\) y \(1728\) de la gramática: cuentan objetos de tipos distintos.
La tabla externa se abre después de fijar identidad, fibra, registro,
operador, proyector y coordenatización.

La reconstrucción de una razón positiva mediante una torre que recibe
esa razón demuestra capacidad representativa del sistema de evaluación.
La salida independiente de una especie exige, en cambio, que los
coeficientes de la torre procedan de su ruta y de su frontera. Esta
separación conserva tanto la completitud representativa como los
resultados sectoriales efectivos. El conjunto queda organizado por una
regla común: construir la estructura y su memoria, aplicar el lector
del tipo correspondiente y efectuar después el reconocimiento físico.
